%% ──────────────────────────────────────────────
%%  پیوست ب – نقد: «مست گیرند، در شهر هر آنکه هست گیرند»
%%  فایل: appendices/app-b-critical-essay.tex
%% ──────────────────────────────────────────────
\chapter{نقد: گر رسم شود که مست گیرند}
\label{ch:critical-essay}

\begin{tcolorbox}[mybox, title={درباره‌ی این پیوست}]
این پیوست — بر خلاف متن اصلی کتاب که تلاش می‌کند بی‌طرفانه و تحلیلی باشد — نقدی شخصی و صریح است. صدای نویسنده در اینجا بلندتر و نقادانه‌تر از فصل‌های اصلی خواهد بود. موضوع: \textbf{چرا مشروعیت‌بخشی به مداخله‌ی خارجی — حتی با بهترین نیت‌ها — باز کردن جعبه‌ی پاندوراست.}
\end{tcolorbox}

%% ================================================================
\section{مسئله: وقتی «حق مداخله» قانونی شود}
\label{sec:app-b-problem}
%% ================================================================

\needspace{6\baselineskip}
فرض کنید پذیرفته‌ایم که در شرایطی استثنایی — نسل‌کشی، جنایات جنگی گسترده — مداخله‌ی خارجی مشروع و حتی ضروری است. این فرض معقول به نظر می‌رسد. اما مسئله اینجاست: \textbf{چه کسی تعریف می‌کند «شرایط استثنایی» چیست؟}

در یک نظام بین‌المللی بدون حکومت مرکزی — آنچه نظریه‌پردازان \concept{«آنارشی بین‌المللی»} می‌نامند — قدرتمندترین بازیگران تعریف‌کنندگان واقعی هنجارها هستند. وقتی «حق مداخله» به رسمیت شناخته شود:

\begin{itemize}
\item آمریکا از آن برای حمله به عراق استفاده می‌کند (۲۰۰۳)
\item روسیه از آن برای الحاق کریمه استفاده می‌کند (۲۰۱۴): «حمایت از روس‌زبانان»
\item عربستان از آن برای بمباران یمن استفاده می‌کند (۲۰۱۵): «بازگرداندن دولت مشروع»
\item ترکیه از آن برای حمله به شمال سوریه استفاده می‌کند: «مبارزه با تروریسم»
\item اسرائیل از آن برای هر عملیاتی استفاده می‌کند: «دفاع مشروع»
\end{itemize}

«مست گیرند، در شهر هر آنکه هست گیرند» — وقتی مشروبخانه باز شود، فقط مشروبخوران محترم نمی‌آیند.

\begin{tcolorbox}[enemybox, title={تز اصلی این پیوست}]
\textbf{مشروعیت‌بخشی اصولی به مداخله‌ی خارجی، جعبه‌ی پاندورایی باز می‌کند که بستنش غیرممکن است.} زیرا:
\begin{enumerate}
\item هیچ مکانیزم بی‌طرفی برای تشخیص «مداخله‌ی مشروع» از «تجاوز مشروعیت‌یافته» وجود ندارد
\item قدرتمندان همیشه می‌توانند بحران را تعریف (یا ایجاد) کنند
\item هر سابقه‌ی «مداخله‌ی موفق» مجوزی برای مداخلات بعدی (و معمولاً بدتر) صادر می‌کند
\item ملت‌های ضعیف قربانیان اصلی این «حق» هستند — نه بهره‌برداران آن
\end{enumerate}
\end{tcolorbox}


%% ================================================================
\section{نقد R2P: ابزار نجات یا ابزار سلطه؟}
\label{sec:app-b-r2p-critique}
%% ================================================================

\needspace{6\baselineskip}
دکترین «مسئولیت حمایت» (R2P) با بهترین نیت‌ها تدوین شد — واکنشی به رواندا ۱۹۹۴. اما یک دهه پس از تصویب رسمی‌اش (۲۰۰۵)، کارنامه‌اش چنین است:

\begin{enumerate}
\item \textbf{لیبی ۲۰۱۱:} تنها مورد استفاده‌ی کامل R2P → نتیجه: هرج‌ومرج، بازار برده، بحران مهاجرت — \textbf{R2P کار خودش را کرد، اما نتیجه فاجعه بود}
\item \textbf{سوریه ۲۰۱۱–:} R2P اعمال نشد (وتوی روسیه و چین) → نتیجه: ۵۰۰٬۰۰۰ کشته — \textbf{R2P نتوانست کار خودش را بکند}
\item \textbf{یمن ۲۰۱۵–:} R2P حتی مطرح نشد (عربستان متحد غرب) → نتیجه: بدترین بحران انسانی — \textbf{R2P نخواست کار خودش را بکند}
\item \textbf{میانمار ۲۰۱۷/۲۰۲۱:} R2P بی‌اثر بود → نسل‌کشی روهینگیا و کودتا — \textbf{R2P نمی‌تواند کار خودش را بکند}
\end{enumerate}

نتیجه: R2P در تنها موردی که اعمال شد (لیبی) فاجعه آفرید، و در موارد دیگر یا نخواست یا نتوانست عمل کند. این دقیقاً مشکل «حق مداخله» است: \textbf{یا سوءاستفاده می‌شود، یا بی‌اثر می‌ماند. حالت سومی وجود ندارد.}


%% ================================================================
\section{نقد نجات‌طلبان: آیا ملتی حق دارد بیگانه را دعوت کند؟}
\label{sec:app-b-invitee-critique}
%% ================================================================

%\needspace{6\baselineskip}
بخشی از ادبیات نجات‌طلبی بر «دعوت» تأکید دارد: «ما خودمان از بیگانه دعوت کرده‌ایم — پس مشکل کجاست؟» اما:

\begin{tcolorbox}[wavebox, title={پنج اشکال در استدلال «دعوت»}]
\textbf{اشکال اول: چه کسی دعوت می‌کند؟}
معمولاً «دعوت‌کننده» یک فرد، حزب، یا گروه خاص است — نه کل ملت. هادی در یمن، کرزی در افغانستان، شورای انتقالی در لیبی — هیچ‌کدام نمایندگان واقعی «ملت» نبودند.

\textbf{اشکال دوم: آیا دعوت آزادانه است؟}
در شرایط بحران، ناامیدی، و ترس — دقیقاً شرایطی که فصل ۱۱ تحلیل کرد — «انتخاب آزاد» معنای واقعی خود را از دست می‌دهد. فردی که غرق می‌شود به هر ریسمانی چنگ می‌زند — اما این به معنای «انتخاب آزاد» نیست.

\textbf{اشکال سوم: آیا دعوت قابل پس‌گرفتن است؟}
وقتی نیروی خارجی وارد شد، «دعوت‌کننده» دیگر کنترلی بر آن ندارد. عراقی‌هایی که از سقوط صدام خوشحال بودند، کنترلی بر سیاست‌های برمر نداشتند. افغان‌هایی که از سقوط طالبان استقبال کردند، کنترلی بر بمباران‌های هوایی نداشتند.

\textbf{اشکال چهارم: آیا نسل‌های بعد هم «دعوت» کرده‌اند؟}
مداخله‌ی خارجی اثرات بلندمدتی دارد که نسل‌های آینده را درگیر می‌کند — نسل‌هایی که هرگز «دعوت» نکرده‌اند.

\textbf{اشکال پنجم: آیا «دعوت» ساختگی نیست؟}
تاریخ نشان می‌دهد که بسیاری از «دعوت‌ها» ساختگی بودند: «دعوت» دولت افغانستان از شوروی (۱۹۷۹) در واقع سازمان‌دهی‌شده‌ی خود شوروی بود. «درخواست کمک» اپوزیسیون لیبی عمدتاً از سوی گروه‌های تبعیدی هماهنگ با غرب بود.
\end{tcolorbox}


%% ================================================================
\section{دفاع از حاکمیت: نه دفاع از دیکتاتوری}
\label{sec:app-b-sovereignty}
%% ================================================================

\needspace{6\baselineskip}
منتقدان ممکن است بگویند: «دفاع از حاکمیت = دفاع از دیکتاتوری.» این استدلال نادرست و خطرناک است.

دفاع از \concept{اصل حاکمیت} به معنای دفاع از \textit{هر رژیمی} نیست — به معنای دفاع از حق مردم هر کشور برای تعیین سرنوشت خودشان بدون دخالت بیگانه است. تفاوت بنیادین:

\begin{itemize}
\item \textbf{نقض حاکمیت از بیرون:} مداخله‌ی نظامی، کودتا، تغییر رژیم — آنچه این کتاب نقد کرد
\item \textbf{تغییر از درون:} انقلاب، اصلاحات، مقاومت مدنی — آنچه این کتاب به‌عنوان جایگزین پیشنهاد کرد
\end{itemize}

\thinker{جان استوارت میل} (\lat{John Stuart Mill}) — خود یک لیبرال تمام‌عیار — در ۱۸۵۹ نوشت: «اگر مردمی آزادی‌شان را خودشان به دست نیاورده باشند، آزادی‌شان را حفظ نخواهند کرد.» این نه دفاع از دیکتاتوری، بلکه دفاع از \concept{عاملیت مردم} است.


%% ================================================================
\section{نقد «واقع‌گرایانه» نجات‌طلبی}
\label{sec:app-b-realist-critique}
%% ================================================================

\needspace{6\baselineskip}
حتی اگر استدلال اخلاقی (حق حاکمیت) را نپذیریم، استدلال \concept{واقع‌گرایانه} (\lat{Realist}) علیه نجات‌طلبی همچنان قوی است:

\begin{tcolorbox}[failbox, title={پنج دلیل واقع‌گرایانه علیه نجات‌طلبی}]
\begin{enumerate}
\item \textbf{آمار علیه نجات‌طلبی است:} از ده‌ها مورد بررسی‌شده، کمتر از ۱۰٪ نتیجه‌ی «نسبتاً مثبت» داشته‌اند. این بدترین «نرخ موفقیت» ممکن برای هر سیاستی است.

\item \textbf{هزینه‌ی فرصت عظیم است:} ۸ تریلیون دلار هزینه‌ی جنگ‌های پس از ۱۱ سپتامبر می‌توانست آموزش ابتدایی برای همه‌ی کودکان جهان را برای ۳۰۰ سال تأمین کند.

\item \textbf{اثر بومرنگ واقعی است:} مداخله تقریباً همیشه مشکلات جدیدی ایجاد می‌کند: عراق → داعش، لیبی → بحران ساحل، افغانستان → بی‌ثباتی پاکستان.

\item \textbf{پیش‌بینی‌ناپذیری:} حتی بهترین تحلیل‌گران نمی‌توانند نتایج بلندمدت مداخله را پیش‌بینی کنند. سیستم‌های اجتماعی-سیاسی پیچیده‌اند و رفتار غیرخطی دارند.

\item \textbf{مشوق‌های نهادی:} مجتمع نظامی-صنعتی، شرکت‌های امنیتی خصوصی، و سیاستمداران جنگ‌طلب انگیزه‌ی مادی برای ادامه‌ی مداخلات دارند — حتی وقتی شکست آشکار است.
\end{enumerate}
\end{tcolorbox}


%% ================================================================
\section{جمع‌بندی: «آزادی‌ای که از بیرون بیاید...»}
\label{sec:app-b-conclusion}
%% ================================================================

\needspace{6\baselineskip}
این نقد ادعا نمی‌کند که جهان باید در برابر نسل‌کشی و جنایت سکوت کند. ادعایش این است:

\begin{enumerate}
\item \textbf{ابزارهای غیرنظامی} (دیپلماسی، مذاکره، تحریم هدفمند، دادگاه‌های بین‌المللی) باید اولویت داشته باشند
\item \textbf{مداخله‌ی نظامی} فقط — و فقط — در شرایطی بسیار خاص و با شروط سخت‌گیرانه مشروع است
\item \textbf{«حق مداخله»} نباید به اصل تبدیل شود — باید استثنای استثناها باشد
\item \textbf{عاملیت مردم} مهم‌ترین عامل تغییر پایدار است — و هر سیاستی که عاملیت مردم را تضعیف کند (حتی به نام «نجات»)، در نهایت زیان‌بار است
\end{enumerate}

\begin{tcolorbox}[quotebox, title={سخن پایانی}]
«حکمت هزاران سال تاریخ بشر به ما می‌گوید: هیچ‌کس به‌اندازه‌ی خود ملت‌ها انگیزه‌ی واقعی برای نجاتشان ندارد — و هیچ‌کس به‌اندازه‌ی قدرت‌های خارجی انگیزه‌ی واقعی برای سوءاستفاده از نام نجات ندارد. میان این دو واقعیت، راه باریکی وجود دارد: راه عاملیت، خوداتکایی، و تفکر انتقادی. این کتاب تلاشی بود — هر قدر ناقص — برای روشن کردن آن راه.»\\
— نویسنده
\end{tcolorbox}

\cleardoublepage